\documentclass[a4paper,11pt]{ltjsarticle}
\usepackage{amsmath,amssymb,amsthm,mathtools}
\usepackage[hidelinks]{hyperref}

\theoremstyle{plain}
\newtheorem{theorem}{定理}[section]
\newtheorem{proposition}[theorem]{命題}
\newtheorem{lemma}[theorem]{補題}
\newtheorem{conjecture}[theorem]{予想}
\theoremstyle{definition}
\newtheorem{definition}[theorem]{定義}
\newtheorem{remark}[theorem]{注意}
\newtheorem{problem}[theorem]{問題}

\newcommand{\N}{\mathbb{N}}
\newcommand{\Z}{\mathbb{Z}}
\newcommand{\Q}{\mathbb{Q}}
\newcommand{\R}{\mathbb{R}}
\newcommand{\dd}{\,\mathrm{d}}

\title{オイラーの定数の無理性について\\
{\large ガンマ関数による表示の厳密化と，無理性の証明に必要な条件}}
\author{Masaaki Yamaguchi}
\date{2026年11月16日}

\begin{document}
\maketitle

\begin{abstract}
オイラーの定数 $\gamma$ を，ゼータ関数の極限およびガンマ関数の微分係数
$\gamma=-\Gamma'(1)=-\int_0^\infty e^{-x}\log x\dd x$ として表す二つの表示を，
完全な証明とともに与える．
続いて，これらの表示から $\gamma$ の無理性を導こうとする議論を検討し，
積分表示だけでは数の算術的性質は決まらないこと，
および無理性の証明には整数係数の一次形式 $q_n\gamma-p_n$ の構成が必要十分であることを示す．
$\gamma$ の無理性は本稿執筆時点で未解決であり，本稿はこれを予想として述べ，
証明のために埋めるべき箇所を問題として定式化する．
\end{abstract}

\section{序論}

オイラーの定数（オイラー・マスケローニ定数）は
\[
  \gamma=\lim_{n\to\infty}\Bigl(\sum_{k=1}^{n}\frac1k-\log n\Bigr)=0.5772156649\ldots
\]
で定義される．$\gamma$ が無理数であるかどうかは，
オイラー以来の未解決問題である \cite{Lagarias}．

著者はノート \cite{Note} において，$\gamma$ を
\[
  C=\int\frac{1}{x^{s}}\dd x-\log x
\]
と書き，ガンマ関数 $\Gamma$ とその導関数
$\Gamma'=\int e^{-x}x^{t-1}\log x\dd x$ を経由して
「オイラーの定数はガンマ関数から無理数と証明される」と述べた．
本稿の目的は，この議論のうち厳密に成立する部分を定理として確定し，
成立しない部分を明示して，無理性の主張を正確な形に置き直すことである．

結論を先に述べる．
\begin{itemize}
  \item ノートの出発点である二つの表示（ゼータ関数による表示とガンマ関数による表示）は正しく，
        定理~\ref{thm:zeta}，定理~\ref{thm:gamma} として証明できる．
  \item これらの表示から無理性を導く部分は証明になっていない（第~\ref{sec:review}~節）．
        したがって本稿は $\gamma\notin\Q$ を定理ではなく予想~\ref{conj:irr} として述べる．
  \item 無理性の証明に必要十分な条件は命題~\ref{prop:criterion} で与えられ，
        ガンマ関数の表示からこの条件を満たす列を構成することが残された問題である
        （問題~\ref{prob:main}）．
\end{itemize}

\section{オイラーの定数の定義と存在}

\begin{definition}
$H_n=\sum_{k=1}^{n}1/k$ とおき，$\gamma=\lim_{n\to\infty}(H_n-\log n)$ と定める．
\end{definition}

実数 $s\in[1,2]$ と自然数 $n\ge1$ に対して
\[
  a_n(s)=\frac{1}{n^{s}}-\int_n^{n+1}\frac{\dd x}{x^{s}}
        =\int_n^{n+1}\Bigl(\frac{1}{n^{s}}-\frac{1}{x^{s}}\Bigr)\dd x
\]
とおく．

\begin{lemma}\label{lem:bound}
$s\in[1,2]$，$n\ge1$ に対して $0\le a_n(s)\le 2/n^{2}$ が成り立つ．
\end{lemma}

\begin{proof}
$x\in[n,n+1]$ のとき $x^{-s}\le n^{-s}$ であるから $a_n(s)\ge0$．
平均値の定理より，ある $\xi\in[n,x]$ があって
$n^{-s}-x^{-s}=s\,\xi^{-s-1}(x-n)\le s\,n^{-s-1}\le 2n^{-2}$．
これを $[n,n+1]$ 上で積分すればよい．
\end{proof}

\begin{proposition}\label{prop:exist}
極限 $\gamma$ は存在し，$\gamma=\sum_{n=1}^{\infty}a_n(1)$，$0<\gamma<1$ である．
\end{proposition}

\begin{proof}
$a_n(1)=1/n-\log\bigl((n+1)/n\bigr)$ であるから
$\sum_{n=1}^{N}a_n(1)=H_N-\log(N+1)$．
補題~\ref{lem:bound} により左辺は $N\to\infty$ で収束し，
$\log(N+1)-\log N\to0$ であるから $H_N-\log N$ も同じ値に収束する．
各 $a_n(1)$ は正であるから $\gamma>0$．
また $n\ge1$ に対して $\log(1+1/n)>1/(n+1)$ であるから
$a_n(1)<1/n-1/(n+1)$ となり，$\gamma<\sum_{n\ge1}\bigl(1/n-1/(n+1)\bigr)=1$．
\end{proof}

\section{ゼータ関数による表示}

ノート \cite{Note} の式 $C=\int x^{-s}\dd x-\log x$ は，
次の定理として正確に述べることができる．
$s>1$ に対して $\zeta(s)=\sum_{n\ge1}n^{-s}$ とする．

\begin{theorem}\label{thm:zeta}
\[
  \gamma=\lim_{s\to1+0}\Bigl(\zeta(s)-\frac{1}{s-1}\Bigr).
\]
\end{theorem}

\begin{proof}
$s>1$ のとき $\int_1^\infty x^{-s}\dd x=1/(s-1)$ であるから
\[
  \zeta(s)-\frac{1}{s-1}=\sum_{n=1}^{\infty}a_n(s)\qquad(1<s\le2).
\]
補題~\ref{lem:bound} とワイエルシュトラスの優級数判定法により，
右辺の級数は $[1,2]$ 上で一様収束する．
各 $a_n$ は $[1,2]$ 上連続であるから，和も $[1,2]$ 上連続である．
よって $s\to1+0$ の極限は $\sum_{n}a_n(1)$ に等しく，
命題~\ref{prop:exist} によりこれは $\gamma$ である．
\end{proof}

\section{ガンマ関数による表示}

$t>0$ に対して
\[
  \Gamma(t)=\int_0^\infty e^{-x}x^{t-1}\dd x
\]
とする．積分記号下の微分により
$\Gamma'(t)=\int_0^\infty e^{-x}x^{t-1}\log x\dd x$ である
（$t$ がコンパクト区間 $[t_0,t_1]\subset(0,\infty)$ を動くとき，
被積分関数の $t$ 微分は可積分関数
$e^{-x}(x^{t_0-1}+x^{t_1-1})\lvert\log x\rvert$ で押さえられる）．

\begin{lemma}\label{lem:beta}
$n\ge0$ に対して
\[
  \int_0^1(1-u)^{n}\log u\dd u=-\frac{H_{n+1}}{n+1}.
\]
\end{lemma}

\begin{proof}
$v=1-u$ と置換し，$\log(1-v)=-\sum_{k\ge1}v^{k}/k$ $(0\le v<1)$ を用いる．
各項は同符号であるから単調収束定理により項別積分ができて
\[
  \int_0^1 v^{n}\log(1-v)\dd v
  =-\sum_{k=1}^{\infty}\frac{1}{k(n+k+1)}
  =-\frac{1}{n+1}\sum_{k=1}^{\infty}\Bigl(\frac1k-\frac{1}{n+k+1}\Bigr)
  =-\frac{H_{n+1}}{n+1}.
\]
最後の等号は，部分和が $H_K-(H_{n+K+1}-H_{n+1})$ であり，
$H_{n+K+1}-H_K\to0$ $(K\to\infty)$ となることによる．
\end{proof}

\begin{remark}
補題~\ref{lem:beta} の左辺は，ベータ関数
$B(a,n+1)=\int_0^1u^{a-1}(1-u)^{n}\dd u$ の $a$ に関する微分係数の $a=1$ での値である．
ノート \cite{Note} の表題にあるベータ関数は，この形で $\gamma$ の表示に現れる．
\end{remark}

\begin{theorem}\label{thm:gamma}
\[
  \gamma=-\Gamma'(1)=-\int_0^\infty e^{-x}\log x\dd x.
\]
\end{theorem}

\begin{proof}
$f_n(x)=(1-x/n)^{n}\log x$ $(0<x\le n)$，$f_n(x)=0$ $(x>n)$ とおく．
$0\le(1-x/n)^{n}\le e^{-x}$ $(0\le x\le n)$ であり，
$n\to\infty$ で $(1-x/n)^{n}\to e^{-x}$ であるから，
$\lvert f_n(x)\rvert\le e^{-x}\lvert\log x\rvert$ かつ $f_n(x)\to e^{-x}\log x$．
$e^{-x}\lvert\log x\rvert$ は $(0,\infty)$ 上可積分であるから，ルベーグの優収束定理により
\[
  \int_0^\infty e^{-x}\log x\dd x=\lim_{n\to\infty}\int_0^{n}\Bigl(1-\frac{x}{n}\Bigr)^{n}\log x\dd x .
\]
$x=nu$ と置換し，$\int_0^1(1-u)^{n}\dd u=1/(n+1)$ と補題~\ref{lem:beta} を用いると
\[
  \int_0^{n}\Bigl(1-\frac{x}{n}\Bigr)^{n}\log x\dd x
  =n\int_0^1(1-u)^{n}(\log n+\log u)\dd u
  =\frac{n}{n+1}\bigl(\log n-H_{n+1}\bigr).
\]
$n\to\infty$ のとき $n/(n+1)\to1$，$H_{n+1}-\log n\to\gamma$ であるから，右辺は $-\gamma$ に収束する．
\end{proof}

\begin{remark}
ノート \cite{Note} では $\Gamma=\int e^{-x}x^{1-t}\dd x$，$\gamma=\Gamma'$ と書かれているが，
指数は $t-1$，符号は $\gamma=-\Gamma'(1)$ が正しい．
\end{remark}

\section{無理性の判定条件}

\begin{proposition}\label{prop:criterion}
実数 $\alpha$ について，次の二条件は同値である．
\begin{enumerate}
  \item $\alpha$ は無理数である．
  \item 整数列 $(p_n)$，$(q_n)$ で，すべての $n$ について $q_n\alpha-p_n\ne0$ であり，
        かつ $q_n\alpha-p_n\to0$ $(n\to\infty)$ となるものが存在する．
\end{enumerate}
\end{proposition}

\begin{proof}
(2)$\Rightarrow$(1)．$\alpha=a/b$（$a\in\Z$，$b\in\N$）と仮定する．
$q_n\alpha-p_n=(q_na-p_nb)/b$ の分子は $0$ でない整数であるから
$\lvert q_n\alpha-p_n\rvert\ge1/b$ となり，$0$ への収束に反する．

(1)$\Rightarrow$(2)．ディリクレの近似定理により，各 $n\ge1$ に対して
$1\le q_n\le n$ なる整数 $q_n$ と整数 $p_n$ で $\lvert q_n\alpha-p_n\rvert<1/n$ となるものがある．
$\alpha$ は無理数であるから $q_n\alpha-p_n\ne0$ である．
\end{proof}

したがって，$\gamma$ の無理性を証明することは，
条件 (2) を満たす整数列を $\alpha=\gamma$ に対して具体的に構成することと同値である．
$\zeta(3)$ の無理性に関するアペリの証明や，$e$，$\pi$ の無理性の古典的な証明は，
すべてこの形をしている．

\section{ノートの議論の検討}\label{sec:review}

ノート \cite{Note} の議論のうち，定理~\ref{thm:zeta} と定理~\ref{thm:gamma} に対応する部分は
前節までで厳密に示された．それ以降の部分については，以下の点が証明として成立していない．

\begin{enumerate}
  \item \textbf{記号の定義がない．}
        $\dd\mathrm{vol}$，$\dd x_m$，$\dd m$ による「単体積分」「多様体積分」は，
        積分領域と測度が指定されておらず，値が定まらない．
  \item \textbf{定数による微分．}
        $\dfrac{\dd}{\dd f}$，$\dfrac{\dd}{\dd\gamma}$ は $f=\gamma$ という定数に関する微分であり，
        意味をもたない．
  \item \textbf{定数と関数の等置．}
        最終式 $2\cos(ix\log x)$ は $x$ の関数であるが，$\gamma$ は定数である．
        また恒等式としては $e^{-f}-e^{f}=-2\sinh f=2i\sin(if)$ であり，
        $2\cos(if)=e^{f}+e^{-f}$ とは一致しない．
  \item \textbf{算術的な議論がない．}
        どの段階にも，整数・有理数の性質（分母の評価など）を用いる箇所がない．
        命題~\ref{prop:criterion} が示すとおり，無理性は整数との関係についての性質であり，
        解析的な等式の変形だけからは得られない．
\end{enumerate}

特に (4) は本質的である．ガンマ関数の積分で表されるという事実は，
数が無理数であることを含意しない．実際，
\[
  \Gamma(n+1)=\int_0^\infty e^{-x}x^{n}\dd x=n!\in\Z,\qquad
  \Gamma\bigl(\tfrac12\bigr)=\int_0^\infty e^{-x}x^{-1/2}\dd x=\sqrt{\pi}\notin\Q
\]
のように，同じ形の積分が有理数にも無理数にもなる．
したがって「ガンマ関数で表されるから無理数である」という推論は成り立たず，
$\gamma=-\Gamma'(1)$ に固有の算術的情報を取り出す追加の議論が必要である．

\section{既知の結果}

$\gamma$ の算術的性質について現在知られている主な結果は次のとおりである．
いずれも無理性そのものには到達していない．

\begin{itemize}
  \item 連分数展開の数値計算により，$\gamma$ が有理数 $a/b$ であるならば
        分母 $b$ は $10^{242080}$ より大きい（Papanikolaou，\cite{Lagarias} 参照）．
  \item Mahler \cite{Mahler} は，ベッセル関数 $J_0$，$Y_0$ に対して
        $\dfrac{\pi}{2}\dfrac{Y_0(2)}{J_0(2)}-\gamma$ が超越数であることを示した．
  \item Sondow \cite{Sondow} は，二重積分
        $\displaystyle\iint_{[0,1]^2}\frac{-\bigl(x(1-x)y(1-y)\bigr)^{n}}{(1-xy)\log xy}\dd x\dd y$
        から $\gamma$ と対数を含む一次形式を構成し，
        $\gamma$ の無理性の十分条件を整数列の言葉で与えた．
  \item Aptekarev \cite{Aptekarev} は，$\gamma$ とオイラー・ゴンペルツ定数
        $\delta=\int_0^\infty\dfrac{e^{-x}}{1+x}\dd x$ の少なくとも一方が無理数であることを示し，
        Rivoal \cite{Rivoal} はこれを「少なくとも一方が超越数である」に強めた．
  \item Murty--Saradha \cite{MurtySaradha} は，オイラー・レーマー定数 $\gamma(a,q)$ の族において
        代数的数は高々一つであることを示した．
\end{itemize}

\section{結論と残された問題}

\begin{conjecture}\label{conj:irr}
オイラーの定数 $\gamma$ は無理数である．
\end{conjecture}

本稿で証明されたのは，$\gamma$ の二つの表示（定理~\ref{thm:zeta}，定理~\ref{thm:gamma}）と，
無理性の判定条件（命題~\ref{prop:criterion}）である．
予想~\ref{conj:irr} は証明されていない．
ノート \cite{Note} の方針，すなわちガンマ関数の表示から無理性を導くという方針を
証明として完成させるには，次の問題を解く必要がある．

\begin{problem}\label{prob:main}
表示 $\gamma=-\int_0^\infty e^{-x}\log x\dd x$ から出発して，
整数列 $(p_n)$，$(q_n)$ と積分の列 $I_n$ で
\[
  I_n=q_n\gamma-p_n,\qquad I_n\ne0,\qquad I_n\to0\quad(n\to\infty)
\]
を満たすものを構成せよ．
\end{problem}

定理~\ref{thm:gamma} の証明に現れた積分
$\int_0^{n}(1-x/n)^{n}\log x\dd x=\dfrac{n}{n+1}(\log n-H_{n+1})$
は $\gamma$ の近似を与えるが，$\log n$ が残るため $q_n\gamma-p_n$ の形にはなっていない．
対数の項を消去すること，あるいは対数を含む一次形式のままで無理性を結論できる
評価を与えること（Sondow \cite{Sondow} の方法）が，問題~\ref{prob:main} の核心である．

\begin{thebibliography}{9}
\bibitem{Note}
M.~Yamaguchi,
「オイラーの定数とガンマ関数，そして微分幾何の量子化と虚数による逆三角関数，ベータ関数による場の理論」，
未公刊ノート．

\bibitem{Lagarias}
J.~C.~Lagarias,
Euler's constant: Euler's work and modern developments,
\textit{Bull. Amer. Math. Soc. (N.S.)} \textbf{50} (2013), 527--628.

\bibitem{Mahler}
K.~Mahler,
Applications of a theorem by A.~B.~Shidlovski,
\textit{Proc. Roy. Soc. London Ser. A} \textbf{305} (1968), 149--173.

\bibitem{Sondow}
J.~Sondow,
Criteria for irrationality of Euler's constant,
\textit{Proc. Amer. Math. Soc.} \textbf{131} (2003), 3335--3344.

\bibitem{Aptekarev}
A.~I.~Aptekarev,
On linear forms containing the Euler constant,
preprint, arXiv:0902.1768 (2009).

\bibitem{Rivoal}
T.~Rivoal,
On the arithmetic nature of the values of the gamma function, Euler's constant, and Gompertz's constant,
\textit{Michigan Math. J.} \textbf{61} (2012), 239--254.

\bibitem{MurtySaradha}
M.~R.~Murty and N.~Saradha,
Euler--Lehmer constants and a conjecture of Erd\H{o}s,
\textit{J. Number Theory} \textbf{130} (2010), 2671--2682.
\end{thebibliography}

\end{document}
